%% kapa: exam=1 round=23 year=2012 period=1 form=A subject=civil_law subject_ko=민법 seq=34 number=34
%% source: Q-Net 감정평가사 기출문제 (cst00309) artlSeq=5052578 file=raw/1008/5052578/2012년 제23회 감정평가사 제1차시험 기출문제(1교시 A형).pdf
%% conversion: kordoc 4.24.51 -> tools/gen1.py (auto); review=proofread-sonnet; flags=none
\begin{question}{34}
저당권에 관한 설명으로 옳지 않은 것은? (다툼이 있으면 판례에 의함)
\begin{choices}
  \item 장래에 발생할 특정의 조건부 채권을 피담보채권으로 하는 근저당권을 설정하는 것도 가능하다.
  \item 구분건물의 전유부분에 설정된 저당권의 효력은 특별한 사정이 없는 한 전유부분의 소유자가 나중에 취득한 대지사용권에 미친다.
  \item 저당권자는 채권 전부의 변제를 받을 때까지 저당부동산 전부에 관하여 그 저당권을 행사할 수 있다.
  \item 저당권자는 저당물의 공용징수로 인하여 저당권설정자가 받을 금전의 지급 전에 이를 압류함으로써 저당권을 행사할 수 있다.
  \item 공유자 중 1인의 지분에 저당권이 설정된 후 공유토지가 분할된 경우, 저당권은 저당권설정자가 분할받은 토지에만 효력이 미친다.
\end{choices}
\end{question}
\answer{5}
